\section*{Hoofdstuk \ref{ch:g11:drift-to-spreading} --- Van continentverschuiving tot oceaanbodemspreiding}

\begin{solution}{exo:g11:drift-to-spreading:1}
\omterm{def:g11:drift-to-spreading:spreading}{Oceaanbodemspreiding} is het ontstaan van nieuwe oceaanbodem bij de kam van een
mid-oceanische rug, waar gesmolten gesteente opstijgt en stolt, gevolgd door de
beweging ervan weg van de rug, naar beide kanten. De volledige snelheid is de
snelheid waarmee de twee platen uit elkaar gaan; de halve snelheid is de snelheid
waarmee de bodem aan één kant zich van de as verwijdert. Bij symmetrische
spreiding is volledige snelheid $= 2
\times$ halve snelheid.
\end{solution}

\begin{solution}{exo:g11:drift-to-spreading:2}
Boven zijn \omterm{def:g11:drift-to-spreading:remanence}{Curietemperatuur}, ongeveer \qty{580}{\celsius}, kan magnetiet geen
eigen magnetisatie behouden; pas wanneer het daaronder afkoelt, worden zijn
korrels langs het veld gemagnetiseerd en behouden ze die richting. Opnieuw verhit
tot \qty{700}{\celsius}, boven de \omterm{def:g11:drift-to-spreading:remanence}{Curietemperatuur}, verliest het gesteente zijn
oude verslag; bij het opnieuw afkoelen legt het het veld van die nieuwe tijd vast.
\end{solution}

\begin{solution}{exo:g11:drift-to-spreading:3}
Het geluid gaat heen en terug: $d = 1500 \times 5.2/2 = \qty{3900}{m}$.
\end{solution}

\begin{solution}{exo:g11:drift-to-spreading:4}
Het huidige normale interval begon \num{0.78} miljoen jaar geleden: het duurt al
\num{780000} jaar. De Jaramillo-gebeurtenis duurde $1.07 - 0.99 = \num{0.08}$
miljoen jaar, \num{80000} jaar. Als je de grenzen telt die jonger zijn dan vijf
miljoen jaar (0.78, 0.99, 1.07, 1.77, 1.95, 2.14, 2.15, 2.581, 3.04, 3.11,
3.22, 3.33, 3.58, 4.18, 4.29, 4.48, 4.62, 4.80, 4.89, 4.98): twintig ompolingen.
\end{solution}

\begin{solution}{exo:g11:drift-to-spreading:5}
$65\,000/40\,000 \approx 1.6$: de ruggen zouden ongeveer anderhalve keer rond de
Aarde gaan.
\end{solution}

\begin{solution}{exo:g11:drift-to-spreading:6}
$20/0.78 \approx \qty{25.6}{km}$ per miljoen jaar, dat is een halve snelheid van
ongeveer \qty{2.6}{cm} per jaar; de volledige snelheid is ongeveer \qty{5.1}{cm}
per jaar.
\end{solution}

\begin{solution}{exo:g11:drift-to-spreading:7}
West: $45/3.58 \approx 12.6$~km per miljoen jaar, \qty{1.26}{cm} per jaar;
oost: $44/3.58 \approx 12.3$~km per miljoen jaar, \qty{1.23}{cm} per jaar.
Volledige snelheid: $1.26 + 1.23 \approx \qty{2.5}{cm}$ per jaar. De twee halve
snelheden verschillen ongeveer \qty{2}{\percent}: de spreiding is vrijwel
symmetrisch.
\end{solution}

\begin{solution}{exo:g11:drift-to-spreading:8}
\qty{1.9}{cm} per jaar is \qty{19}{km} per miljoen jaar: $990/19 \approx 52$
miljoen jaar, iets ouder dan de 49 miljoen jaar van het oudste sediment
(ongeveer \qty{6}{\percent}).
\end{solution}

\begin{solution}{exo:g11:drift-to-spreading:9}
Bij \qty{20}{km} per miljoen jaar: \qty{15.6}{km}, \qty{19.8}{km},
\qty{21.4}{km}, \qty{35.4}{km} en \qty{39.0}{km}. Vanaf de as: normaal van 0
tot \qty{15.6}{km}, omgekeerd tot \qty{19.8}{km}, normaal tot \qty{21.4}{km}
(Jaramillo), omgekeerd tot \qty{35.4}{km}, normaal tot \qty{39.0}{km} (Olduvai),
daarna omgekeerd.
\end{solution}

\begin{solution}{exo:g11:drift-to-spreading:10}
$64 \times 10^6/4.6 \times 10^9 \approx \num{0.014}$, ongeveer
\qty{1.4}{\percent} van de ouderdom van de Aarde. Het ondersteunt het argument
van Hess dat de oceaanbodem jong is omdat hij voortdurend wordt gemaakt en
vernietigd: van het grootste deel van de geschiedenis van de Aarde is geen
oceaanbodem overgebleven.
\end{solution}

\begin{solution}{exo:g11:drift-to-spreading:11}
Op elk moment ontstaat de bodem bij de as met de polariteit van dat moment, en
splitst hij: de ene helft verwijdert zich naar de ene kant, de andere naar de
andere, met dezelfde snelheid. Twee stroken die op hetzelfde moment ontstonden,
liggen daarom op dezelfde afstand van de as, één aan elke kant, met dezelfde
polariteit: elke kant is een kopie van de andere, gespiegeld in de as. De
symmetrie wordt verbroken als de bodem aan de ene kant sneller wegschuift dan aan
de andere (asymmetrische spreiding), of als de as naar een nieuwe plaats
verspringt.
\end{solution}

\begin{solution}{exo:g11:drift-to-spreading:12}
(1) \qty{5}{cm} is $5 \times 10^{-5}$~\unit{km}: elk jaar komt er een strook van
$65\,000 \times 5 \times 10^{-5} = \qty{3.25}{km^2}$ bij.
(2) $2.9 \times 10^8/3.25 \approx 9 \times 10^7$ jaar, ongeveer 90 miljoen
jaar. (3) Dat is van dezelfde orde als de ouderdom van de oudste bodem
(190 miljoen jaar) en korter: om de oceaanbodem zo jong te houden, moet er ergens
evenveel bodem worden vernietigd als er bij de ruggen wordt gemaakt.
\end{solution}

\begin{solution}{exo:g11:drift-to-spreading:13}
De gebeurtenis duurde \num{0.08} miljoen jaar. Bij \qty{12.5}{km} per miljoen
jaar is haar streep $0.08 \times 12.5 = \qty{1}{km}$ breed; bij \qty{75}{km} per
miljoen jaar $0.08 \times 75 = \qty{6}{km}$. Gezien vanaf het zeeoppervlak,
\qty{3}{km} hoger, geeft een streep die veel smaller is dan haar diepte maar een
zwak veld, dat vervaagt in de velden van haar brede buren met de tegengestelde
polariteit; een streep die twee keer zo breed is als haar diepte, valt op.
\end{solution}

\begin{solution}{exo:g11:drift-to-spreading:14}
Plaats 16: $176/11 \approx \qty{16}{m}$ per miljoen jaar; plaats 20: $72/67
\approx \qty{1.1}{m}$ per miljoen jaar. Het sediment hoopt zich van plaats tot
plaats en van tijd tot tijd met heel verschillende snelheden op, dus zijn dikte
geeft niet de ouderdom van de bodem. De ouderdom wordt gegeven door de fossielen
van het oudste sediment, dat direct op het basalt ligt.
\end{solution}

\begin{solution}{exo:g11:drift-to-spreading:15}
(1) West: $58/2.581 \approx \qty{22.5}{km}$ per miljoen jaar,
\qty{2.25}{cm} per jaar; oost: $70/2.581 \approx \qty{27.1}{km}$ per miljoen
jaar, \qty{2.71}{cm} per jaar; volledige snelheid $128/2.581 \approx \qty{49.6}{km}$ per
miljoen jaar, ongeveer \qty{5.0}{cm} per jaar. (2) Ongeveer \qty{225}{km} ten
westen en \qty{271}{km} ten oosten. (3) De twee grenzen lagen samen bij de as
toen de ompoling plaatsvond; sindsdien zijn ze met de volledige snelheid uit
elkaar gegaan, hoe de beweging ook over de twee kanten is verdeeld. Hun afstand
gedeeld door hun ouderdom is de volledige snelheid.
\end{solution}

\begin{solution}{pb:g11:drift-to-spreading:1}
\textbf{1.} Plaats 16 ligt het dichtstbij (\qty{191}{km}, 11 miljoen jaar),
plaats 20 het verst (\qty{1270}{km}, 67 miljoen jaar): hoe verder van de rug,
hoe ouder de bodem.

\textbf{2.} Het sediment begint zich op het basalt af te zetten zodra het basalt
is ontstaan, dus de oudste laag is ongeveer even oud als het basalt. Ze kan alleen
iets jonger zijn, nooit ouder.

\textbf{3.} In kilometer per miljoen jaar: plaats 16: $191/11 \approx 17.4$;
plaats 15: $380/24 \approx 15.8$; plaats 18: $506/26 \approx 19.5$; plaats 17:
$643/33 \approx 19.5$; plaats 14: $727/40 \approx 18.2$; plaats 19: $990/49
\approx 20.2$; plaats 20: $1270/67 \approx 19.0$.

\textbf{4.} $129.5/7 \approx \num{18.5}$~km per miljoen jaar, ongeveer
\qty{1.85}{cm} per jaar.

\textbf{5.} Ja: hun verhoudingen (19.5) liggen tussen die van de westelijke
plaatsen. De bodem is aan beide kanten met ongeveer dezelfde snelheid gespreid:
de spreiding is symmetrisch.

\textbf{6.} $t = d/u$: de ouderdom is evenredig met de afstand.

\textbf{7.} $u = 208\,276/10\,952 \approx \num{19.0}$~km per miljoen jaar,
dat is \qty{1.90}{cm} per jaar.

\textbf{8.} Ongeveer $2 \times 1.9 = \qty{3.8}{cm}$ per jaar.

\textbf{9.} Ja: \qty{1.9}{cm} per jaar is iets minder dan \qty{2}{cm} per
jaar.

\textbf{10.} $1270/19.0 \approx 67$ miljoen jaar, evenveel als het oudste
sediment ervan.

\textbf{11.} $191/9 \approx \num{21.2}$~km per miljoen jaar,
\qty{2.1}{cm} per jaar.

\textbf{12.} $1270/71 \approx \num{17.9}$~km per miljoen jaar,
\qty{1.8}{cm} per jaar.

\textbf{13.} $180\,329/9854 \approx \num{18.3}$~km per miljoen jaar,
\qty{1.83}{cm} per jaar.

\textbf{14.} $(1.90 - 1.83)/1.83 \approx \qty{4}{\percent}$: de twee snelheden
komen overeen binnen de onzekerheden van de ouderdommen en afstanden.

\textbf{15.} De ene ouderdom komt van fossielen in het sediment, de andere van het
magnetisme van het basalt en de tijdschaal van de ompolingen. Een fout in de ene
methode zou in de andere niet verschijnen: hun overeenstemming is een sterke toets
van de \omterm{def:g11:drift-to-spreading:spreading}{oceaanbodemspreiding}.

\textbf{16.} $13.8 + 13.8 = \ang{27.6}$.

\textbf{17.} $111.2 \times \cos 30^\circ \approx \qty{96.3}{km}$ per graad;
$27.6 \times 96.3 \approx \qty{2660}{km}$.

\textbf{18.} $2660/19.0 \approx 140$ miljoen jaar.

\textbf{19.} $2660/130 \approx \num{20.5}$~km per miljoen jaar, ongeveer
\qty{2.0}{cm} per jaar, iets meer dan \qty{1.9}{cm} per jaar. Omdat de laatste
67 miljoen jaar ongeveer \qty{1.9}{cm} per jaar haalden, ging de oceaan in het
begin iets sneller open: de bodem tussen plaats 20 en de rand, $2660 - 1270 =
\qty{1390}{km}$, ontstond in $130 - 67 = 63$ miljoen jaar, ongeveer
\qty{2.2}{cm} per jaar.

\textbf{20.} \textbf{De oudste oceaanbodem bij Afrika op \ang{30} ZB is ongeveer
130 miljoen jaar oud} (140 miljoen jaar als je de huidige snelheid doortrekt): de
zuidelijke Atlantische Oceaan begon daar in het Krijt open te gaan.
\end{solution}
